\section{4교시 (90분) — 실습 ① 컷오버 런북·SAC + 운영 이관·SLA 3층}

% ===== 4교시 (90분) — 실습 ① 컷오버 런북·SAC + 운영 이관·SLA 3층 =====
\begin{frame}{이 교시에서 하는 것 — 오전 이론 둘을 손으로 쓴다}
\begin{itemize}
\item \textbf{⑯ 런북·SAC} — go/no-go 3 · PoNR · 롤백 트리거 4 · SAC 10항목을 \textbf{직접} 쓴다
\item \textbf{⑰ 운영 이관} — SLA/OLA/UC 3층 매핑 7행 + error budget 환전을 \textbf{손계산}한다
\item 두 문서는 서로의 입력 — SAC 항목 7(SLA 서명)이 ⑰ 항목 4, 런북 D+1 서명이 ⑰ 항목 1
\item 도구는 종이 캔버스 · 계산기 · 워크북 §1.4·§2.4 빈 템플릿 — 노트북은 검산용
\item 시간 — 안내 10분 / ⑯ 35분 / ⑰ 30분 / 짝 교환 검산 15분
\end{itemize}
\keymsg{정답은 워크북 §1.5·§2.5에 있다 — 보지 말고 쓰고, 다 쓴 뒤에 대조하라}
\note{워크북 §1·§2 실습. Day3 4교시와 같은 진행 — 과제 지시 두 프레임, 작성 요령 네 프레임, 흔한 오류 한 프레임, 예시본 대조, 정리. 오전 2교시(런북 구조·게이트·PoNR·롤백 비대칭)와 3교시(SAC·ELS·SLA 3층·error budget)가 이 90분의 재료다. 두 과제를 하나의 교시에 묶은 이유 — 런북의 마지막 행(D+1 07:20 인수 서명)이 운영 이관의 첫 행이고, SAC 항목 7(SLA/OLA/UC 서명본)은 ⑰ 항목 4의 표가 없으면 판정할 수 없다. 오픈은 이벤트가 아니라 이관이라는 부제를 두 문서가 손으로 증명하게 된다. 예상 소요 3분.}
\end{frame}

\begin{frame}{과제 ⑯ — 온담 1차 컷오버, 창 02-26 20:00\til{}02-28 06:00(34h)을 런북으로}
\begin{itemize}
\item 게이트 3개 — 착수(19:30) · 데이터 정본 전환(01:00) · 매장 배포(12:00) — 조건·판정자·재실행 허용 횟수
\item PoNR 시각과 근거, 기술적 롤백 최종 시한 — \textbf{두 시각을 분리}해 적고 그 사이의 결정권자를 적는다
\item 롤백 트리거 4 — 관측 지표 · 임계값(리허설 실측에서) · 관측 시각 · 판정자 · 전체/부분 롤백
\item SAC 10항목 — 조건(관측 가능) · 판정자(인수자 실명) · 증빙 · G3 사전 점검 · G4 판정
\item 주어진 것 — 리허설 3회 실측(이력 적재 6h 50m) · 매장 620점 · 인터페이스 47본 · PG 7본 · 헌장 §14 완료 기준 ①
\end{itemize}
\keymsg{묻는 것은 하나 — 새벽 3시에 이 종이만 보고 "지금 돌리면 되는가"에 답할 수 있는가}
\note{워크북 §1.4 빈 템플릿 항목 3·4·5·8을 쓴다. 항목 1·2(문서 정보·시각표 38행)와 항목 6·7·9·10은 예시본으로 읽기만 한다 — 시각표를 다 쓰면 35분이 끝난다. 학생이 자주 묻는 것 — "게이트는 왜 세 개인가?" 되돌리는 비용이 0 / 시간 / 시간+손실+매장으로 달라지는 지점이 셋이기 때문이다(2교시 그림 2-2). "SAC 10항목이 워크북의 12항목과 다른가?" 워크북 예시본은 12항목이고 과제는 10항목 이상이면 된다 — 운영 조직의 관심 순서(문서·모니터링·백업복구·보안·교육·KEDB·SLA·절차·소스·인터페이스·계약·잔여 결함)를 지키는지를 본다. 임계값은 반드시 숫자로 — "이상 없음"은 채점하지 않는다. 예상 소요 4분(안내), 작성 35분.}
\end{frame}

\begin{frame}{과제 ⑰ — SLA 7항목을 3층으로 받치고, 99.5\%를 분과 릴리스 횟수로 환전한다}
\begin{itemize}
\item 3층 매핑표 — 행 = SLA 7항목(가용성·P1·P2 복구·배치·발주 마감·출고 확정·데이터 정정)
\item 열 = SLA 목표 / OLA(사내 팀·목표) / UC(외부 계약·목표) / 합산(가장 긴 사슬) / 판정
\item 주어진 UC 초안 — 수행사 유지보수 착수 1h·복구 4h · CSP 99.9\% · 3PL 대성로지스 \textbf{UC 없음}
\item error budget — 99.5\%/월 \ra{} 허용 다운(분) \ra{} 비계획 60\% / 릴리스 40\% \ra{} 릴리스당 기대 다운 \ra{} 월 허용 횟수
\item 주어진 DORA 실측(2026-09) — 변경 실패율 12\% · 평균 복구 240분 / ELS 목표 10\% · 45분
\end{itemize}
\keymsg{묻는 것은 둘 — 어느 행이 구조적으로 불가능한가, 복구 시간을 얼마로 줄여야 주 1회 릴리스가 SLA와 양립하는가}
\note{워크북 §2.4 빈 템플릿 항목 3·4를 쓴다. 항목 2(운영 문서 12종)·5(ELS 8주)·6\til{}9는 예시본으로 읽는다. 3층 매핑은 7행 중 ② P1 복구와 ⑥ 출고 확정 두 행이 연습의 핵심 — 나머지 다섯은 성립하는 행이고 10분 안에 끝낸다. error budget은 계산기로 — 평균 달 30.4375일을 쓰라고 미리 말한다(30일로 계산한 조는 216분이 나오고 그것은 틀린 것이 아니라 다른 전제다). 학생이 자주 묻는 것 — "계획 정지는 다운인가?" SLA 산정 제외, 단 CAB 승인 건만 — 이 규칙 한 줄을 표에 적었는지 본다. 예상 소요 3분(안내), 작성 30분.}
\end{frame}

\begin{frame}{작성 요령 ① — PoNR은 "무엇이 되돌릴 수 없게 되는가"로 정한다}
\fig[0.58]{../그림/PM_Day4/fig_2_2_reversal_cost.pdf}
\figcap{그림 2-2. 되돌리는 비용의 계단 — 게이트 1 전 0 \ra{} 시간만 \ra{} 시간+손실(PoNR 01:40) \ra{} +매장 방문 2.1억(게이트 3) \ra{} 불가(16:00). 계단이 오르는 자리마다 결정권자가 한 층 올라간다}
\keymsg{PoNR = 신 원장이 정본이 되어 이후 거래가 그곳에만 쌓이기 시작하는 시각 — 배포 시작이 아니다}
\note{교재 §2.3, 워크북 ⑯ 항목 4. 온담은 게이트 2 통과(01:38) 직후 신 재고·주문 원장 활성화 01:40이 PoNR — 03:00 3PL 폴백 배치 1회차가 신 원장 기준으로 출고 지시를 만들기 때문에, 그 뒤의 롤백은 수작업 대사 1,200건 이상·2영업일을 남긴다. 기술적 롤백 최종 시한 16:00은 별개의 계산(오픈 06:00 $-$ 복원 8h $-$ 검증 2h $-$ 매장 캐시 재동기화 4h)이다. 두 시각을 하나로 적는 조가 매 기수 절반이다 — PoNR 이후 16:00까지는 "가능하지만 손실 있음"이고 결정권자가 P1 PM에서 정미란으로 올라간다는 것을 표에 적게 한다. 매장 배포를 PoNR로 적은 조에는 "매장은 되돌릴 수 있다, 원장 분기는 되돌리기 어렵다"를 되묻는다. 예상 소요 6분.}
\end{frame}

\begin{frame}{작성 요령 ② — SAC의 판정자는 인수자다, 프로젝트가 판정하면 자기 평가다}
\fig[0.58]{../그림/PM_Day4/fig_3_2_sac_progress.pdf}
\figcap{그림 3-2. SAC 12항목 × 2시점 — G3 사전 점검 01-28(충족 2·미해소 7·대기 3) \ra{} 14주 해소 \ra{} G4 05-21(충족 11·조건부 1). 조건부 1건은 해소 기한·책임자 서명으로 이관 1번이 된다}
\keymsg{헌장 §14 완료 기준 ① — "SAC 미해소 0건 또는 해소 기한·책임자 서명" — 이 한 줄이 판정 규칙이다}
\note{교재 §3.2, 워크북 ⑯ 항목 8. 판정자 열을 먼저 채우게 한다 — 원칙은 인수자 박준영, 기술 항목(소스·인터페이스)은 박세연, 보안은 최영주, 계약 이관은 강현도. 판정자 열에 P1 PM이나 수행사 PM이 들어간 조는 그 행을 다시 쓰게 한다. 증빙 열에 "완료 보고"라고 쓴 조도 마찬가지 — 문서 번호·시험 기록·서명본이어야 한다. G3 사전 점검 열이 있어야 G4에서 처음 보는 일이 없다는 것, 그리고 조건부 1건(항목 11 아카이브 접근 절차, 재경 승인 06-15)이 완료 기준 ①의 두 번째 절로 통과해 ⑱ 이관 1번이 된다는 것이 이 그림의 요점이다. 예상 소요 5분.}
\end{frame}

\begin{frame}{작성 요령 ③ — 3층은 아래에서 위로: UC 없는 SLA는 약속이 아니다}
\fig[0.58]{../그림/PM_Day4/fig_3_4_sla_chain.pdf}
\figcap{그림 3-4. SLA ② P1 복구 2시간 — 초안 사슬 15+30+30+240 = 315분 > 120 / 재협상 165 / 병렬 착수로 $\le$ 120. 항목 ⑥ 출고 확정은 3PL UC 부재로 "최선 노력" 강등}
\keymsg{약속은 그것을 받치는 약속의 합보다 강할 수 없다 — 합산은 가장 긴 사슬로, 한 층이 비면 "구조적 불가"}
\note{교재 §3.5, 워크북 ⑰ 항목 4. 쓰는 순서를 뒤집게 한다 — SLA 목표를 먼저 적고 OLA·UC를 끼워 맞추는 것이 아니라, UC(수행사 복구 4h, CSP 99.9\%, 3PL 없음)와 OLA(착수 30분·원인 분리 30분)를 먼저 적고 합을 올려 SLA가 성립하는지 판정한다. ② P1 복구는 초안 UC로 315분이라 구조적으로 불가 — 선택지는 UC 재협상(착수 30·복구 90, 24/7 대기, 연 +0.6억 [추가 설정])이거나 SLA 문안 강등이다. ⑥ 출고 확정 D+1 06:00은 3PL 대성로지스 UC가 물류본부 계약이라 P1 밖 — "최선 노력"으로 강등하고 UC는 "배치 파일 수신 확인 30분"으로 한정 서명(한동석 05-10). 판정 열에 "성립"만 있고 "구조적 불가"가 하나도 없는 조는 UC를 읽지 않은 것이다. 예상 소요 6분.}
\end{frame}

\begin{frame}{작성 요령 ④ — error budget 산식: 99.5\%를 분으로, 분을 릴리스 횟수로}
\begin{tbl}[\scriptsize]{L{0.30\textwidth}L{0.34\textwidth}L{0.28\textwidth}}
\textbf{단계} & \textbf{산식} & \textbf{온담 값} \\ \midrule
월 허용 다운 & 30.4375일 × 24 × 60 × 0.5\% & \textbf{219.2분} (3h 39m) \\
비계획 유보 / 릴리스 몫 & 60\% / 40\% & 131.5분 / \textbf{87.7분} \\
릴리스당 기대 다운(실측) & 변경 실패율 12\% × 복구 240분 & 28.8분 \ra{} 월 \textbf{3.0회} \\
릴리스당 기대 다운(중간) & 10\% × 90분 & 9.0분 \ra{} 월 9.7회 \\
릴리스당 기대 다운(ELS 목표) & 10\% × 45분 & 4.5분 \ra{} 월 \textbf{19.5회} \\
소진 시 규칙 & 유보 초과 시 표준 릴리스 동결 · 다음 달 초기화 & 결정자 박준영 \\
\end{tbl}
\keymsg{복구 시간을 4시간에서 45분으로 줄이지 않으면 주 1회 릴리스도 SLA와 양립하지 않는다 — 회귀 자동화·KEDB의 존재 이유}
\note{교재 §3.6, 워크북 ⑰ 항목 3. 검산값은 src/day4\_closing.py 출력과 같다. 계산 순서가 곧 전제다 — 평균 달(30.4375일)을 쓰는지, 계획 정지(2차 컷오버 8h·월 1회 유지보수 창)를 SLA 산정에서 제외하는 규칙을 적었는지, 릴리스 몫을 몇 \%로 두는지. 릴리스당 기대 다운 = 변경 실패율 × 평균 복구 시간이라는 한 줄이 CAB의 "안정 vs 기능" 감정 토론을 산수로 바꾼다. 3월 ELS 실측 57분(유보의 43\%, 예산 소진 26\%)을 표 아래 적어 예산이 "쓰이는" 것임을 보인다. 그림 3-5(도넛)는 노트판에서 보조로 쓴다. 예상 소요 6분.}
\end{frame}

\begin{frame}{흔한 오류 — 이 실습에서 매 기수 나오는 다섯}
\begin{enumerate}
\item 게이트가 하나(오픈 직전)뿐이고 조건이 "이상 없음" — 되돌리는 비용이 달라지는 지점이 게이트다
\item PoNR을 "배포 시작"으로 적고 기술적 시한과 한 시각으로 합친다 — 사이의 결정권자가 사라진다
\item 롤백 트리거가 "중대 장애 발생 시" — 임계값 없는 트리거는 새벽에 토론이 된다
\item SAC 판정자가 프로젝트 쪽, 증빙이 "완료 보고" — 인수가 아니라 자기 평가
\item SLA만 있고 OLA는 구두, 3PL·앱 벤더는 UC에서 뺀다 — 99.5\%를 분으로 환산하지 않는다
\end{enumerate}
\keymsg{다섯 오류의 공통점 — 판정자와 숫자를 적지 않았다. 판정자 없는 조건은 조건이 아니다}
\note{워크북 §1.3·§2.3 각 항목의 "흔한 오류" 줄을 모은 것. 짝 교환 검산 15분에서 이 다섯을 체크리스트로 준다 — 상대 조의 런북에서 ① 게이트 수와 조건의 관측 가능성 ② PoNR과 기술적 시한이 다른 시각인가 ③ 트리거 4개 모두에 숫자 임계값이 있는가 ④ SAC 판정자 열에 프로젝트 사람이 있는가 ⑤ 3층 표에 "구조적 불가" 행이 있는가·219.2분이 나왔는가. 세 개 이상 걸리는 조는 예시본 대조 전에 그 항목만 다시 쓴다. 예상 소요 4분.}
\end{frame}

\begin{frame}{예시본 참조 · 시간 배분 · 심화 과제}
\begin{itemize}
\item 예시본 — 워크북 §1.5(런북 38행·게이트 3·트리거 4·SAC 12) · §2.5(문서 12종·3층 7행·error budget·ELS 8주)
\item 인쇄용 — \textsf{../템플릿/PM\_Day4/16\_컷오버런북\_SAC\_완성예시.pdf} · \textsf{17\_운영이관\_SLA\_완성예시.pdf} (빈 양식은 \textsf{\_빈템플릿.pdf})
\item 시간 — ⑯ 게이트·PoNR·트리거·SAC 35분 / ⑰ 3층 매핑·환전 30분 / 짝 교환 검산·대조 15분
\item 자기 점검 — 워크북 §1.7·§2.7 점검 질문 각 3개. 값이 다르면 \textbf{어느 전제가 달랐는지}를 적는다
\item 심화(3\til{}7년차) — 자기 회사의 최근 오픈 당일 기록에서 게이트·PoNR·임계값·편차 기록 넷을 찾아 본다
\end{itemize}
\keymsg{예시본은 정답이 아니라 전제가 적힌 답이다 — 다른 값이 나왔으면 전제를 비교하라}
\note{Day3 4교시와 같은 운영. 예시본 PDF는 실습 시작 시 배포하지 않고 35분 뒤(⑰ 착수 시점)에 ⑯ 예시본만 먼저 배포, ⑰ 예시본은 손계산이 끝난 뒤. 워크북 §1.7 기본 과제(게이트 2에서 재실행 2회 후에도 14건 불일치 \ra{} RB-1 발동 \ra{} 항목 3·4·5·7 재작성, R-14 확률 갱신, 롤백 런북 10행)와 §2.7 기본 과제(3PL 재협상 타결 시 3층 표 재작성·릴리스 몫 40\% 유지 가능성 계산)는 시간이 남는 조의 추가 과제이고, 대개 과제로 낸다. 심화 과제의 네 질문 중 셋 이상이 "아니오"인 조직은 다음 오픈에서 이번 오픈으로부터 아무것도 배우지 못한 채 시작한다는 문장을 그대로 읽어 준다. 예상 소요 90분 전체.}
\end{frame}

\begin{frame}{4교시 핵심 정리}
\begin{enumerate}
\item 게이트 3개는 되돌리는 비용이 달라지는 세 지점 — 조건은 관측값, 판정자는 실명, 재실행 횟수까지
\item PoNR 01:40과 기술적 시한 16:00은 다른 시각 — 그 사이는 "가능하지만 손실 있음", 결정권자가 올라간다
\item 트리거 4개의 임계값은 리허설 실측에서 온다 — 발동 0이 결과이지 목표가 아니다
\item SAC는 인수자가 판정해야 인수다 — 조건부 1건은 기한·책임자 서명으로 이관 목록에 넘어간다
\item 3층은 아래에서 위로, 99.5\%는 219.2분 \ra{} 87.7분 \ra{} 릴리스 3회 — 복구 시간이 릴리스 속도를 정한다
\end{enumerate}
\keymsg{5교시는 종료 — 보고서는 선언이 아니라 증명이고, 미해결 사항은 소유자 실명이 없으면 소멸한다}
\note{조별 산출물 회수. 3층 매핑표 1장은 6교시 시작 전까지 벽에 붙여 두고 갤러리워크로 서로 대조하게 한다 — "구조적 불가" 행의 수가 조마다 다르면 그것이 5교시 미해결 사항 이관(어느 약속이 프로젝트 밖으로 넘어가는가)의 도입 재료다. 예상 소요 3분.}
\end{frame}
